\section{Conclusion}
\label{sec:conclusion}

A qubit that returns to a known state is a gift to a compiler. It is also a
moment at which the system decides who may use that physical resource next.
Those two facts should not be separated.

Coherent teleportation shows the constructive side. Deferred corrections factor
the source and ancilla into $\ket{++}$, and two Hadamards return them to
$\ket{00}$. The postcondition is exact, compositional, and suitable for reuse.

The spy-detector analysis shows the adversarial side at two complementary
levels. The complete randomized export restores the baseline computational-
basis probabilities on the displayed register while lowering purity to $1/2$
and leaving a distinguishable hidden branch. The fixed-input reconstruction then
identifies that branch: after the ordered SWAP network, $q_4$ stores the spy's
measurement result. It is a perfect record of Alice's value when the spy uses
Alice's basis and is maximally mixed when the bases differ. At the same time,
the reconstructed victim subsystem remains equal to its clean routed reference
to floating-point precision.

The general theorem explains the boundary. Orthogonal, classical information
can be copied into hidden workspace without changing its source. Nonorthogonal
pure states cannot reveal distinguishable information to an environment while
remaining unchanged. Between those cases lies the engineering problem: users
often verify only a subset of observables, compilers operate across abstraction
layers, and real programs contain both quantum coherence and accessible
classical structure.


The capability matrix turns that boundary into a reviewer-navigable security
map. It shows, without treating color as evidence, which objectives are directly
feasible, which require probabilistic success or accepted disturbance, which
are plausible but unvalidated, and which are blocked under their stated
assumptions. Its starred cluster connects the project-specific source-access
result to retained measurement routing and delayed use of side information,
while the linked registry records what each mechanism touches and what it may
leave unchanged.

The public version-0.6.0 artifact makes the result reproducible through exact
branch enumeration, numerical validation, tests across multiple Python
versions, and deterministic output comparison in continuous integration. That
software evidence narrows the remaining questions rather than erasing them. The
next decisive task is to reconstruct the original detector's actual Boolean
acceptance register, followed by independent Qiskit parity, formal
information-flow contracts, and controlled compiler and hardware experiments.

Secure quantum compilation should therefore treat ``clean'' as more than a
single-qubit value. A released qubit should have a proven state, proven
separability, explicit liveness, and an information-flow guarantee. The same
analysis that makes scarce qubits reusable can help ensure that they are not
reused against the computation they were meant to serve.
